% =====================================================================
% Section 3 — Method  (revised)
% Drop-in replacement for the current Sec. 3. Requires amsmath, amssymb.
% Decision points flagged with %% NOTE: ... — toggle as needed.
% =====================================================================

\section{Method}
\label{sec:method}

\subsection{Overview}
\label{sec:overview}
The basis of our design is simple: within one image, manipulated and
pristine patches are \emph{relatively} different, and a strong self-supervised
representation already exposes that difference once it is lightly adapted. A
frozen DINOv3 backbone, adapted only by low-rank (LoRA) updates, produces dense
patch features (Sec.~\ref{sec:backbone}) that are fed to two lightweight
\emph{linear} heads: a within-image patch-contrastive head for localization and
a gated multiple-instance (MIL) pooling head for image-level prediction. The
heads are deliberately light, serving primarily to insulate the backbone and to
aggregate its features, respectively, and to inject the two training signals
that shape the shared backbone. At inference, localization is performed by
$k$-means clustering the output of the contrastive projection head and then
resolving which cluster is the manipulation using the MIL pooling weights
(Sec.~\ref{sec:decode}); there is no separate decoder and no per-pixel
probabilities. A training-free, attention-guided zoom re-runs the same model on
the suspected region to stabilize and sharpen the result.

\subsection{Frozen backbone with two lightweight heads}
\label{sec:backbone}
We use a DINOv3 ViT-H\textsuperscript{+}/16 backbone~\cite{dinov3} at input
resolution $448\times448$
%% NOTE: 448 is the default operating point (swept in Sec. 5.2); the deliverable
%% checkpoint trains at higher resolution. Keep/adjust to match the headline.
, which tokenizes the image into a $28\times28$ grid of $N{=}784$ patch tokens.
The backbone weights are frozen; we inject LoRA adapters~\cite{lora}
(rank $r{=}16$, $\alpha{=}32$, dropout $0.1$) into the attention and MLP
projection layers, so the only trainable parameters are the LoRA updates and the
two heads. From the shared patch features $\mathbf{f}\in\mathbb{R}^{N\times D}$
the model produces:
\begin{enumerate}\itemsep2pt
  \item an $L_2$-normalized contrastive embedding
        $\mathbf{z}_i\in\mathbb{R}^{d}$ ($d{=}64$) per patch, from a
        \emph{single linear} projector; and
  \item an image-level forgery logit together with a set of per-patch pooling
        weights $a_i$, from a gated-attention MIL head in the sense of Ilse
        \etal~\cite{ilse2018mil}.
\end{enumerate}
We stress that $a_i$ are MIL \emph{bag-pooling weights}, not spatial or
transformer attention: the head is two linear gates and a linear scorer, has no
localization objective, and is never supervised on which patches are real or
fake. Both heads are deliberately low-capacity, and---crucially---the
contrastive head projects the $D$-dimensional features down to $d{=}64$ before
the loss is applied. Imposing the contrastive objective on a compact projection,
rather than on the backbone features directly, confines task-specific shaping to
the head and keeps the rich self-supervised features near their semantic roots,
which gradients pulling patches together and apart would otherwise corrupt. The
heads thus serve mainly to route the two training signals into the shared
backbone through LoRA; at inference they are read-outs---the projector emits the
embeddings we cluster, the MIL head emits the polarity weight.

\subsection{Training objective}
\label{sec:objective}
The localization signal comes from a within-image, pairwise contrastive loss on
the patch embeddings. We treat labeled pristine and modified regions
symmetrically to sidestep the issues classical splices raise---both the pasted
region and the host are real pixels from some source, so the label is ill-posed.
Simply put, our contrastive objective asks that same-label patches be alike and
different-label patches be unlike, treating the two labeled regions symmetrically.

For an image with binary patch labels $y_i\in\{0,1\}$ (pristine / manipulated),
let $s_{ij}=\mathbf{z}_i^{\top}\mathbf{z}_j$ be the cosine similarity. We use two
hard thresholds---a cohesion threshold $\tau^{+}{=}0.60$ and a separation
threshold $\tau^{-}{=}0.15$---each enforced by a one-sided hinge that is zero
once its threshold is met (and never pushes cohesive pairs apart):
\begin{equation}
  v^{+}_{ij}=\big[\tau^{+}-s_{ij}\big]_{+}\ (y_i{=}y_j),
  \qquad
  v^{-}_{ij}=\big[s_{ij}-\tau^{-}\big]_{+}\ (y_i{\neq}y_j).
\end{equation}
We average these over the three kinds of pairs in an image---same-manipulated
($mm$), same-pristine ($pp$), and cross ($\times$):
\begin{equation}
  A_{\text{in}}=\overline{v^{+}}(mm),\quad
  A_{\text{out}}=\overline{v^{+}}(pp),\quad
  R=\overline{v^{-}}(\times),
\end{equation}
where the per-set mean is count-tempered,
$\ \overline{v}(P)=\big(\textstyle\sum_{(i,j)\in P} v_{ij}\big)\big/\big(a_P^{\rho}\,t_P^{\,1-\rho}\big)$,
with $t_P$ pairs in $P$, $a_P$ of them violating, and $\rho{=}0.75$ (denominators
floored at $1$). The exponent $\rho$ interpolates between averaging over
violators only ($\rho{=}1$) and over all pairs ($\rho{=}0$), so both the severity
and the \emph{number} of violations move the loss. The per-image loss balances
the two cohesion terms by region area and adds separation:
\begin{equation}
  \ell_{\text{con}}=w_{\text{in}}A_{\text{in}}+w_{\text{out}}A_{\text{out}}+\lambda_{r}R,
  \qquad
  w_{\text{in}}=\frac{p^{\beta}}{p^{\beta}+(1-p)^{\beta}},\ \ w_{\text{out}}{=}1{-}w_{\text{in}},
\end{equation}
with $p$ the manipulated-patch fraction, $\beta{=}0.5$ (a small splice is neither
drowned nor over-amplified), and $\lambda_{r}{=}1$. Authentic images and crops
containing a single region keep only their surviving cohesion term and are
down-weighted by $\kappa{=}0.05$ in the batch mean $\mathcal{L}_{\text{con}}$.

An auxiliary image-level head supplies the second signal. Its gated-MIL pooling
produces an image logit supervised by binary cross-entropy
$\mathcal{L}_{\text{bce}}$ (``is this image manipulated''); the same pooling
weights $a_i$ are reused at decode time for cluster polarity. The total loss is
\begin{equation}
  \mathcal{L}
  = \lambda_{\text{bce}}\,\mathcal{L}_{\text{bce}}
  + \lambda_{\text{con}}\,\mathcal{L}_{\text{con}},
  \qquad \lambda_{\text{bce}}{=}1,\ \lambda_{\text{con}}{=}2 .
\end{equation}
%% NOTE: an optional dense patch-BCE head exists but is omitted from the main
%% model; it is used only to instantiate the supervised-threshold decoder
%% baseline in the ablations (Sec. 5).

\subsection{GT-free decoding and attention-guided zoom}
\label{sec:decode}
At inference a single forward pass yields the patch embeddings, the pooling
weights, and the image logit. We do \emph{not} threshold a trained mask head.
Instead we cluster the embeddings $\{\mathbf{z}_i\}$ into two groups with
spherical $k$-means ($k{=}2$) and take the resulting partition as the localization.

\paragraph{Cluster polarity.}
Clustering alone cannot say \emph{which} of the two groups is the manipulation;
the partition is correct but its foreground/background assignment is
arbitrary~\cite{huh2018fighting}. We resolve this with the MIL pooling weights:
the manipulated label is assigned to the cluster with higher mean $a_i$. This is
worth dwelling on, because it is the one place the system relies on a signal it
never directly trains. The image head has no localization objective and is never
told which patches are manipulated; its pooling weights concentrate on the
manipulated region only as a \emph{by-product} of the image-level real/fake
task---the patches most useful for that decision are usually the edited ones. We
quantify this emergence in Sec.~\ref{sec:ablations}: the pooling map aligns with
the ground-truth manipulation at AUC well above chance, and that alignment
collapses to chance when the contrastive objective is removed, even though the
two heads share no gradient path---evidence that the localization the pooling
map exposes is borrowed from the contrastively-shaped backbone, not learned by
the head. The reliance is honest but imperfect: when the manipulated region
occupies a large fraction of the image the polarity can invert, since nothing
constrains it. This is a known, unsupervised trade-off, and we report where it
occurs.

\paragraph{Attention-guided zoom.}
A single coarse pass is fragile for two reasons, only one of which is grid
resolution. First and foremost, two-way $k$-means is unstable: on small or
low-contrast edits a poor partition can annex large, semantically-unrelated
regions alongside (or instead of) the true edit, producing a mask that is
spuriously spread across the frame. Second, even a correct coarse partition is
imprecise about \emph{what exactly} changed inside a region the model flagged as
suspicious. We therefore add a training-free refinement (``zoom''): from the
coarse pooling map we derive a bounding box around the attended region, crop the
input to that box, re-run the same frozen model on the crop, and re-cluster.
Inside the smaller, more homogeneous crop the partition is far less likely to
wander, the model resolves the edited sub-region more finely, and---as a third
benefit---the effective patch density over the region of interest increases. The
box derivation, padding, and place-back are pure geometry and touch neither
ground truth nor a trained head. We report both the coarse pass (``stage~1'')
and the refined pass (``stage~2'').

We use $k$-means as the default decoder throughout; a density-based alternative
(HDBSCAN) and a supervised-threshold decoder are examined in
Sec.~\ref{sec:ablations}.

\subsection{Leakage-free evaluation}
\label{sec:leakage}
Forgery-localization pipelines are unusually prone to silent ground-truth
leakage, so we enforce three structural invariants, checked by automated tests:
the model forward pass exists in exactly one module; ground-truth masks are read
in exactly one module (the metric computation); and the decode functions are
structurally GT-free---they accept only the model-output bundle, which has no
mask field, so a decoder cannot read ground truth even by accident. Cluster
polarity is therefore chosen by the pooling weights, never by comparison to the
ground-truth mask.

\subsection{Implementation details}
\label{sec:impl}
We adapt DINOv3 ViT-H\textsuperscript{+}/16 (patch size $16$) with LoRA rank
$16$, $\alpha{=}32$, dropout $0.1$ on the attention and MLP projections; the
contrastive projector has dimension $d{=}64$ and the MIL pool hidden dimension
is $256$. We optimize with AdamW (learning rate $2\!\times\!10^{-4}$, weight
decay $10^{-4}$), a one-epoch linear warmup followed by cosine decay, gradient
clipping at $1.0$, and mixed-precision (fp16). Loss weights are
$\lambda_{\text{bce}}{=}1$, $\lambda_{\text{con}}{=}2$; contrastive margins
$\tau^{+}{=}0.60$, $\tau^{-}{=}0.15$, $\lambda_r{=}1$, norm-power $\rho{=}0.75$,
area-balance power $\beta{=}0.5$, single-class weight $\kappa{=}0.05$. Training
runs to best validation localization F1 (minimum $5$ epochs, early-stop patience
$2$); ablation cells are capped at a fixed epoch budget for comparability. At
inference, $k$-means uses $k{=}2$ with $4$ restarts. Models were trained on
NVIDIA L4 / A100 (Colab) and dual RTX~2080~Ti hardware; timing is reported on a
single RTX~2080~Ti.
